\documentclass[11pt]{article}
\usepackage[utf8]{inputenc}
\usepackage{amsmath,amssymb,amsthm}
\usepackage[margin=1in]{geometry}
\usepackage{enumitem}
\usepackage{hyperref}

\begin{document}
% ==================== PAGE 17 ====================
\begin{enumerate}[resume]
    \item $M_n(\mathbb{Z}) \to \text{set of all } n \times n \text{ matrices with integer entries is also a non-commutative ring with unity.}$
    \item $A = \left\{ \begin{bmatrix} x & x \\ x & x \end{bmatrix} \;\middle|\; x \in \mathbb{R} \right\}$ is a commutative ring with unity.
\end{enumerate}

\noindent $(A, +)$ is an abelian group.

\noindent Let $X, Y \in A$ s.t. $X = \begin{bmatrix} x & x \\ x & x \end{bmatrix}$, $Y = \begin{bmatrix} y & y \\ y & y \end{bmatrix}$
\[
\begin{bmatrix} x & x \\ x & x \end{bmatrix} \begin{bmatrix} y & y \\ y & y \end{bmatrix} = \begin{bmatrix} 2xy & 2xy \\ 2xy & 2xy \end{bmatrix} = \begin{bmatrix} y & y \\ y & y \end{bmatrix} \begin{bmatrix} x & x \\ x & x \end{bmatrix}
\]
Let $E$ be identity of $A$, s.t. $E = \begin{bmatrix} e & e \\ e & e \end{bmatrix}$
\[
\begin{bmatrix} x & x \\ x & x \end{bmatrix} \begin{bmatrix} e & e \\ e & e \end{bmatrix} = \begin{bmatrix} x & x \\ x & x \end{bmatrix}
\]
\[
\begin{bmatrix} 2xe & 2xe \\ 2xe & 2xe \end{bmatrix} = \begin{bmatrix} x & x \\ x & x \end{bmatrix}
\]
\[
2xe = x \implies e = \frac{1}{2}
\]
\[
\therefore E = \begin{bmatrix} 1/2 & 1/2 \\ 1/2 & 1/2 \end{bmatrix}
\]

\noindent For making $(A, \cdot)$ abelian group: \\
If $Y = X^{-1}$ then
\[
\begin{bmatrix} x & x \\ x & x \end{bmatrix} \begin{bmatrix} y & y \\ y & y \end{bmatrix} = \begin{bmatrix} 1/2 & 1/2 \\ 1/2 & 1/2 \end{bmatrix}
\]
\[
2xy = \frac{1}{2} \implies y = \frac{1}{4x} \quad (x \ne 0)
\]
\[
\therefore Y = \begin{bmatrix} \frac{1}{4x} & \frac{1}{4x} \\ \frac{1}{4x} & \frac{1}{4x} \end{bmatrix}
\]

\newpage
% ==================== PAGE 18 ====================
\noindent 8) $A = \left\{ \begin{bmatrix} x & 0 \\ 0 & 0 \end{bmatrix} \;\middle|\; x \in \mathbb{R} \right\}$ is a commutative ring with unity. \\
Here $E = \begin{bmatrix} 1 & 0 \\ 0 & 0 \end{bmatrix}$.

\bigskip
\noindent 9) $R = \{0\}$ is a ring with unity $0$, called \textbf{trivial ring} or \textbf{zero ring}. \\
If $R \ne \{0\}$, then $R$ is called \textbf{non-trivial ring}.

\bigskip
\noindent 10) $\mathbb{Z}[x] =$ set of all polynomials with integer coefficients. \\
$(\mathbb{Z}[x], +, \cdot)$ is a commutative ring with unity $1$.

\bigskip
\noindent 11) $\mathbb{R}[x], \mathbb{Q}[x], \mathbb{C}[x]$ are also commutative rings with unity.

\bigskip
\noindent 12) $C[a, b] =$ set of all continuous real-valued functions defined on $[a, b]$. \\
Define:
\[
(f + g)(x) = f(x) + g(x)
\]
\[
(f \cdot g)(x) = f(x) \cdot g(x) \quad \forall\, x \in [a, b]
\]
Then $(C[a, b], +, \cdot)$ is a commutative ring with unity $e(x) = 1, \; \forall\, x \in [a, b]$.

\bigskip
\noindent \textbf{Properties of a Ring:} \\
Let $(R, +, \cdot)$ be a ring. Then $\forall\, a, b, c \in R$:
\begin{enumerate}
    \item $a \cdot 0 = 0 \cdot a = 0$
    \item $a(-b) = (-a)b = -(ab)$
    \item $(-a)(-b) = ab$
    \item $a(b - c) = ab - ac$
    \item $(b - c)a = ba - ca$
\end{enumerate}

\newpage
% ==================== PAGE 19 ====================
\noindent \textbf{Proof of Properties:}
\begin{enumerate}
    \item $a \cdot 0 = a \cdot (0 + 0) = a \cdot 0 + a \cdot 0$ \\
    $\Rightarrow a \cdot 0 + 0 = a \cdot 0 + a \cdot 0$ \\
    $\Rightarrow a \cdot 0 = 0$ (by left cancellation in $(R, +)$). \\
    Similarly, $0 \cdot a = 0$.

    \item $a \cdot (-b) + ab = a(-b + b) = a \cdot 0 = 0$ \\
    $\Rightarrow a(-b) = -(ab)$. \\
    Similarly, $(-a)b = -(ab)$.

    \item $(-a)(-b) = -(a(-b)) = -(-(ab)) = ab$.

    \item $a(b - c) = a(b + (-c)) = ab + a(-c) = ab - ac$.

    \item $(b - c)a = (b + (-c))a = ba + (-c)a = ba - ca$.
\end{enumerate}

\bigskip
\noindent \textbf{Zero Divisors:} \\
Let $R$ be a ring. A non-zero element $a \in R$ is said to be a \textbf{zero divisor} (or divisor of zero) if there exists a non-zero element $b \in R$ such that:
\[
ab = 0 \quad \text{or} \quad ba = 0
\]
\begin{itemize}
    \item If $ab = 0$ with $a \ne 0, b \ne 0$, then $a$ is called a \textbf{left zero divisor} and $b$ is called a \textbf{right zero divisor}.
\end{itemize}

\bigskip
\noindent \textbf{Examples:}
\begin{enumerate}
    \item In $(\mathbb{Z}_6, +_6, \times_6)$, elements $2, 3, 4$ are zero divisors:
    \[
    2 \times_6 3 = 0, \quad 3 \times_6 4 = 0
    \]
    \item In $(\mathbb{Z}, +, \cdot)$, $(\mathbb{Q}, +, \cdot)$, $(\mathbb{R}, +, \cdot)$, there are no zero divisors.
\end{enumerate}

\newpage
% ==================== PAGE 20 ====================
\noindent \textbf{Integral Domain:} \\
A commutative ring $R$ with unity ($1 \ne 0$) is called an \textbf{integral domain} if it has no zero divisors. \\
i.e., $\forall\, a, b \in R$,
\[
ab = 0 \implies a = 0 \quad \text{or} \quad b = 0
\]

\bigskip
\noindent \textbf{Examples of Integral Domains:}
\begin{enumerate}
    \item $(\mathbb{Z}, +, \cdot)$, $(\mathbb{Q}, +, \cdot)$, $(\mathbb{R}, +, \cdot)$, $(\mathbb{C}, +, \cdot)$ are integral domains.
    \item $(\mathbb{Z}_p, +_p, \times_p)$ is an integral domain if and only if $p$ is a prime.
    \item $\mathbb{Z}[x], \mathbb{Q}[x], \mathbb{R}[x]$ are integral domains.
    \item $\mathbb{Z}[i] = \{a + bi \mid a, b \in \mathbb{Z}\}$ (Gaussian integers) is an integral domain.
\end{enumerate}

\bigskip
\noindent \textbf{Non-examples:}
\begin{enumerate}
    \item $(\mathbb{Z}_n, +_n, \times_n)$ where $n$ is composite is \textbf{not} an integral domain.
    \item $M_n(\mathbb{R})$ ($n \ge 2$) is not an integral domain (neither commutative nor free of zero divisors).
\end{enumerate}

\bigskip
\noindent \textbf{Theorem:} \\
A commutative ring $R$ is an integral domain if and only if the cancellation law holds in $R$:
\[
ab = ac \quad (a \ne 0) \implies b = c
\]

\medskip
\noindent \textbf{Proof:} \\
$(\implies)$ Let $R$ be an integral domain and let $ab = ac$ with $a \ne 0$.
\begin{align*}
ab - ac = 0 &\implies a(b - c) = 0  \\
&\implies b - c = 0 \quad (\text{since } a \ne 0 \text{ and } R \text{ has no zero divisors})  \\
&\implies b = c
\end{align*}
$(\impliedby)$ Conversely, suppose the cancellation law holds in $R$. \\
Let $ab = 0$ with $a \ne 0$. \\
Then $ab = a \cdot 0$. \\
Since $a \ne 0$, by cancellation law:
\[
b = 0
\]
Hence, $R$ has no zero divisors, so $R$ is an integral domain.
\end{document}
